% ============================================================
%  Chapter 17: Universal Transport and Partition Extinction
%  Source: bounded-phase-space-trajectory.tex, Consequences 31--33
% ============================================================

\section{The Universal Transport Formula}
\label{sec:ch17-transport}

\begin{consequence}{Universal Transport Formula}{transport}
  Any transport coefficient $\Xi$ (electrical conductivity, thermal
  conductivity, viscosity, diffusivity) of a many-body system satisfying
  BPSL is given by:
  \begin{equation}\label{eq:transport}
    \Xi = \frac{1}{N}\sum_{p}\sum_{i,j} \tau_p^{(ij)}\, g_{ij},
  \end{equation}
  where $\tau_p^{(ij)}$ is the partition lag time between cells $i$ and $j$
  for carrier $p$, $g_{ij}$ is the partition conductance between those cells,
  and $N$ is the number of carriers.
\end{consequence}

Special cases:
\begin{itemize}
  \item \textbf{Drude resistivity}: $\rho = m/(ne^2\tau_p)$ --- set
    $\tau_p^{(ij)} = \tau_p$ (uniform), $g_{ij} = e^2/m$.
  \item \textbf{Chapman--Enskog viscosity}: $\eta = \half mn\bar{v}\lambda$ ---
    set $\tau_p^{(ij)}$ to the mean free time between cell collisions.
  \item \textbf{Einstein diffusivity}: $D = \kB T\mu_\mathrm{mob}$ --- set
    $\tau_p^{(ij)}$ to the drift time between cells.
\end{itemize}

Formula~\eqref{eq:transport} is a single equation that unifies all three.

\section{Partition Extinction and Zero Resistance}
\label{sec:ch17-extinction}

\begin{consequence}{Partition Extinction at $T_c$}{partitionextinction}
  When the temperature falls below a critical value $T_c$, distinct partition
  cells coalesce into a single cell under phase-locking:
  \begin{equation}
    \lim_{T \to T_c^-} \tau_p^{(ij)} \to 0 \quad \text{for all } (i,j).
  \end{equation}
  Transport then vanishes identically: $\Xi = 0$.
\end{consequence}

\begin{remark}
  Partition extinction gives an exact derivation of superconductivity
  (zero electrical resistance) and superfluidity (zero viscosity).
  The $\tau_p \to 0$ limit is partition cell coalescence, not a phenomenological
  condensate.
  The transition temperature $T_c$ is set by the partition cell splitting energy,
  which in metals equals the BCS pairing energy.
\end{remark}

\section{Wiedemann--Franz Law}
\label{sec:ch17-wf}

\begin{consequence}{Wiedemann--Franz Law}{wiedemann}
  The ratio of thermal conductivity $\kappa$ to electrical conductivity
  $\sigma$ satisfies:
  \begin{equation}
    \frac{\kappa}{\sigma T} = \frac{\pi^2}{3}\left(\frac{\kB}{e}\right)^2
    = L_0 = 2.44 \times 10^{-8}\,\mathrm{W\,\Omega\,K^{-2}}.
  \end{equation}
\end{consequence}

\begin{proof}[Proof sketch]
  Both $\kappa$ and $\sigma$ are special cases of the universal transport
  formula~\eqref{eq:transport}.
  The ratio $\kappa/\sigma T$ equals the ratio of the energy-weighted and
  charge-weighted partition lag times, which by the Sommerfeld expansion at
  the Fermi surface gives the Lorenz number $L_0$.
\end{proof}

%%  SOURCE: integrate full derivations from
%%          bounded-phase-space-trajectory.tex, Consequences 31--33.

\bigskip
\begin{tcolorbox}[colback=crownblue!5, colframe=crownblue!60!black,
                  title={\textbf{Chapter 17 Summary}}]
\begin{itemize}[noitemsep]
  \item The universal transport formula $\Xi = \frac{1}{N}\sum_{p,ij}
    \tau_p^{(ij)} g_{ij}$ unifies Drude resistivity, Chapman--Enskog
    viscosity, and Einstein diffusivity.
  \item Partition extinction ($\tau_p \to 0$ at $T_c$) gives exact zero
    resistance and zero viscosity.
  \item The Wiedemann--Franz law is the ratio of two special cases of the
    universal formula.
\end{itemize}
\end{tcolorbox}
